% =============================================
% File: chapters/pl03_thongtin.tex
% Mẫu PL03 — Thông tin kết quả nghiên cứu
% =============================================

\phantomsection
\addcontentsline{toc}{chapter}{Thông tin kết quả nghiên cứu}

{\fontsize{13pt}{15pt}\selectfont

\vspace{0.2cm}
\noindent BỘ GIÁO DỤC VÀ ĐÀO TẠO\\[0.05cm]
\noindent\textbf{TRƯỜNG ĐH CÔNG NGHỆ KỸ THUẬT TP.HCM}

\vspace{0.3cm}
\begin{center}
    {\bfseries THÔNG TIN KẾT QUẢ NGHIÊN CỨU CỦA ĐỀ TÀI}
\end{center}

\vspace{0.1cm}
\noindent\textbf{1. Thông tin chung:}

\vspace{0.1cm}
\noindent- Tên đề tài: \textit{ROBOT TỰ CÂN BẰNG ĐA BẬC TỰ DO}

\vspace{0.15cm}
\noindent- Chủ nhiệm đề tài: \textbf{Trần Minh Cương} \hfill Mã số SV: \underline{23134007}

\vspace{0.1cm}
\noindent- Lớp: \underline{\quad 23134B \quad} \hfill Khoa: \underline{\quad Cơ khí chế tạo máy \quad}

\vspace{0.2cm}
\noindent- Thành viên đề tài:

\vspace{0.1cm}
% Cột Khoa dùng \raggedright để không xuống dòng giữa chữ
% Tổng: 3.0 + 2.0 + 3.8 + 4.7 = 13.5cm
\noindent\begin{tabular}{|p{0.20\linewidth}|p{0.15\linewidth}|p{0.15\linewidth}|p{0.40\linewidth}|}
\hline
\textbf{Họ và tên} & \textbf{MSSV} & \textbf{Khoa} & \textbf{Điện thoại, email} \\
\hline
Nguyễn Anh Hào       & 23134015 & Cơ khí chế tạo máy & 0379667093, 23134015@student.hcmute.edu.vn \\
\hline
Nguyễn Đình Anh Khôi & 21146256 & Cơ khí chế tạo máy & 0903487536, 21146256@student.hcmute.edu.vn \\
\hline
Phùng Mạnh Thắng     & 21146316 & Cơ khí chế tạo máy & 0961992800, 21146316@student.hcmute.edu.vn \\
\hline
Võ Hồng Quân         & 22134012 & Cơ khí chế tạo máy & 0363645485, 22134012@student.hcmute.edu.vn \\
\hline
\end{tabular}

\vspace{0.15cm}
\noindent- Người hướng dẫn: TS. Cái Việt Anh Dũng

\vspace{0.2cm}
\noindent\textbf{2. Mục tiêu đề tài:} Xây dựng hệ thống điều khiển PI kết hợp ANN để điều
khiển thích nghi động cơ DC NF5475 trên vi điều khiển ESP32; đồng thời phân tích động học
cơ cấu 5-bar parallel làm nền tảng cho điều khiển chân robot bipedal wheel.

\vspace{0.2cm}
\noindent\textbf{3. Tính mới và sáng tạo:} Đề tài kết hợp quy trình nhận dạng hàm truyền
động cơ bằng MATLAB (\texttt{tfest}), tối ưu tham số PI bằng giải thuật di truyền (GA) và
xây dựng mạng ANN từ đầu (không dùng framework) để dự đoán $K_p$, $K_i$ thích nghi theo
sai số vận tốc thực, triển khai trực tiếp trên ESP32. Đây là quy trình hoàn chỉnh và có
thể tái hiện, phù hợp với điều kiện nghiên cứu sinh viên.

\vspace{0.2cm}
\noindent\textbf{4. Kết quả nghiên cứu:}
\begin{itemize}[leftmargin=0.5cm, topsep=2pt, itemsep=0pt]
    \item Hàm truyền động cơ NF5475 được nhận dạng với độ chính xác 95,36\%.
    \item Bộ PI tối ưu bằng GA: độ vọt lố $<5\%$, thời gian xác lập $<0{,}5$~s.
    \item Mạng ANN (2 lớp ẩn, 1000 epochs) hội tụ, dự đoán $K_p$, $K_i$ trên 9 mức
    setpoint từ 250 đến 2250 RPM.
    \item Hệ thống ổn định trên ESP32 với chu kỳ điều khiển $\Delta t = 0{,}01$~s.
    \item Phân tích động học thuận/nghịch và Jacobian cơ cấu 5-bar parallel
    ($A_1A_2 = 105$~mm, $L_{act} = 100$~mm, $L_{pas} = 180$~mm).
\end{itemize}

\vspace{0.15cm}
\noindent\textbf{5. Đóng góp về mặt giáo dục và đào tạo, kinh tế -- xã hội, an ninh, quốc
phòng và khả năng áp dụng của đề tài:} Đề tài cung cấp quy trình nghiên cứu hoàn chỉnh từ
lý thuyết đến thực nghiệm về điều khiển thích nghi, có giá trị giáo dục cao cho sinh viên
ngành Kỹ thuật. Kết quả có thể ứng dụng trong phát triển robot phục vụ và hỗ trợ y tế. Chi
phí phần cứng thấp (ESP32, động cơ DC) giúp dễ nhân rộng trong các cơ sở đào tạo.

\vspace{0.15cm}
\noindent\textbf{6. Công bố khoa học của SV từ kết quả nghiên cứu của đề tài}
\textit{(ghi rõ tên tạp chí nếu có) hoặc nhận xét, đánh giá của cơ sở đã áp dụng các kết
quả nghiên cứu (nếu có):} Chưa có công bố / đang chuẩn bị nộp bài.

% ---- Khối ký tên SV — bên phải ----
\vspace{0.5cm}
\begin{flushright}
    Ngày \underline{\quad} tháng \underline{\quad} năm 2026 \\[0.3cm]
    \textbf{SV chịu trách nhiệm chính thực hiện đề tài} \\[0.15cm]
    \textit{(ký, họ và tên)} \\[1.8cm]
\end{flushright}

% ---- Nhận xét người hướng dẫn ----
\vspace{0.1cm}
\noindent\textbf{Nhận xét của người hướng dẫn về những đóng góp khoa học của SV thực hiện đề tài}
\textit{(phần này do người hướng dẫn ghi):}

\vspace{0.2cm}
\noindent\dotfill\\[0.55cm]
\noindent\dotfill

% ---- Khối ký tên Người hướng dẫn — bên phải ----
\vspace{0.5cm}
\begin{flushright}
    Ngày \underline{\quad} tháng \underline{\quad} năm 2026 \\[0.3cm]
    \textbf{Người hướng dẫn} \\[0.15cm]
    \textit{(ký, họ và tên)} \\[1.8cm]
\end{flushright}

}